\section{Introduction}
\label{sec:introduction}

Large language models for code have shown strong results on a wide range of software engineering tasks, including code generation from natural language, cross-language translation, code summarization, bug repair, and program synthesis \citep{wang2021codet5,roziere2023codellama,hui2024qwen25coder,lozhkov2024starcoder2,lu2021codexglue,chen2021evaluating,austin2021program}. These results come mainly from large-scale offline training, in which data for all target tasks is collected in advance, combined into one corpus, and used to train the model over many iterations. This approach works well when the target task distribution is fixed, but it becomes harder to sustain as the number of tasks grows and deployment conditions change.

In long-term deployment, a code assistant cannot remain unchanged after its initial training. New programming languages emerge, existing frameworks change, and the range of capabilities expected of the model keeps growing. Retraining the model from scratch after each change is slow and inflexible, and in many business settings re-using old training data raises privacy concerns that limit approaches based on storing or replaying historical data. A more practical alternative lets the model absorb new tasks one at a time while continuing to handle previously learned tasks without access to old data. Standard supervised learning is not suited to this setting because it assumes the full training set is available at once and that the task distribution is stationary; neither assumption holds when a model is updated incrementally under storage or privacy constraints.

Continual learning studies exactly this problem: a model learns from a stream of tasks without forgetting what it has already learned \citep{parisi2019continual,delange2021continual}. Updates that reduce the loss on the current task tend to overwrite knowledge from earlier tasks, a phenomenon known as catastrophic forgetting \citep{mccloskey1989catastrophic,goodfellow2013empirical}. The code domain makes this problem particularly difficult because tasks vary widely in input format, output format, programming language, and evaluation criterion --- a task sequence may move from English-to-Java generation to Java-to-C\# translation, Ruby summarization, and multilingual executable program synthesis, a level of diversity comparable to the CodeTask-CL benchmark \citep{yadav2023exploring}. Existing continual learning approaches have clear limitations in this setting. Regularization methods \citep{kirkpatrick2017overcoming,zenke2017continual,aljundi2018memory} accumulate importance weights across tasks, but their penalty becomes unreliable on long and structurally diverse sequences. Prompt-based methods \citep{wang2022learning,razdaibiedina2023progressive} keep the backbone frozen but have limited representational capacity for tasks as different as translation, summarization, and repair. Sequential low-rank adaptation methods \citep{wang2023orthogonal,wang2025continual,qian2025treelora} cause destructive interference within a shared low-rank subspace; orthogonality constraints reduce but do not eliminate this problem, and the number of available orthogonal directions shrinks as tasks accumulate.

A further limitation shared by most existing approaches is the assumption that task identity is available at inference time. Methods that drop this assumption typically rely on heuristics such as nearest-centroid assignment, which break down when task embeddings overlap --- a common occurrence in code generation, where problems from different languages can produce nearly identical instruction embeddings. Current evaluation practice compounds these gaps: the most widely used continual learning benchmark for code \citep{yadav2023exploring} covers only four tasks and is scored with text-similarity metrics such as BLEU \citep{papineni2002bleu} and CodeBLEU \citep{ren2020codebleu}, which cannot detect code that looks plausible but fails functionally, and short task sequences hide long-range forgetting effects. Few prior studies combine sequential learning with executable, multilingual, unit-test-based evaluation \citep{chen2021evaluating,cassano2023multipl}.

This paper addresses task-agnostic, parameter-efficient continual learning for sequential code generation. We adopt a routing-based approach that combines per-task expert isolation with a learned probabilistic model for task identification at inference time, avoiding data replay, parameter interference, and the requirement for explicit task labels. The contributions of this paper are as follows.

\begin{enumerate}[leftmargin=*]
\item \textbf{A two-track evaluation platform for continual code generation.} We extend CodeTask-CL \citep{yadav2023exploring} with four additional tasks covering new programming languages and task types, producing a longer, more diverse eight-task sequence scored by text-overlap metrics. We additionally build a nine-language executable benchmark from OSS-Instruct, McEval-Instruct, McEval, MBPP, and MultiPL-E \citep{wei2023magicoder,chai2024mceval,austin2021program,cassano2023multipl}, in which generated programs are evaluated by running them against unit tests.
\item \textbf{A deduplication pipeline for the executable benchmark} that combines exact matching via normalized SHA-256 hashing with a near-duplicate detection stage based on shingling, MinHash, Locality-Sensitive Hashing, and exact Jaccard scoring \citep{jaccard1901distribution,broder1997resemblance,indyk1998approximate}, reducing train--test leakage.
\item \textbf{A systematic empirical comparison} of regularization-based, prompt-based, low-rank-adaptation-based, and online continual learning baselines on sequential programming tasks under a consistent training and evaluation protocol.
\item \textbf{GMM-Router}, a routing-based continual learning framework that trains a task-specific LoRA expert per task on a frozen backbone, fits a diagonal Gaussian Mixture Model by Expectation-Maximization on routing embeddings extracted through a fixed orthonormal projection, and merges expert adapters at the weight level using temperature-scaled softmax routing weights. For executable code generation, the framework further incorporates a traceback-based router that uses runtime failure signals to guide a second inference round.
\end{enumerate}

We evaluate GMM-Router against seven continual learning baselines and a task-identity upper bound on both benchmark tracks. On the diverse CodeTask benchmark, GMM-Router achieves the highest average score (51.62\%) among all task-agnostic methods, approaching the upper bound (54.15\%) with near-zero forgetting. On the executable benchmark, it achieves 34.40\% Pass@1 and the highest Pass@5 (47.75\%) among all methods evaluated.

The remainder of this paper is organized as follows. Section~\ref{sec:background} presents the motivation and the conceptual background needed to follow the proposed method. Section~\ref{sec:data} describes the data collection, formatting, and deduplication pipeline underlying the two benchmark tracks. Section~\ref{sec:methodology} presents the GMM-Router methodology. Section~\ref{sec:evaluation} defines the research questions and the experimental configuration. Section~\ref{sec:results} reports and discusses the experimental results. Section~\ref{sec:related_work} situates GMM-Router within the continual learning and mixture-of-experts literature. Section~\ref{sec:conclusion} concludes the paper and outlines directions for future work.
